\chapter{Literature Review}\label{ch:lit}

VANET authentication schemes optimise different objectives: privacy, efficiency, decentralisation,
accountability or support for mobility. This chapter groups them by approach, notes what each
contributes and what it leaves open for this work, and positions the project.

\section{PKI and Pseudonymous Certificates}

Raya and Hubaux~\cite{raya2007securing} define the VANET threat model and a public-key
infrastructure in which every vehicle carries a large pool of short-lived anonymous certificates.
The authority can resolve a pseudonym when an investigation is authorised. This architecture underlies
the US Security Credential Management System (SCMS)~\cite{brecht2018scms} and IEEE
1609.2~\cite{ieee16092}. Its costs are the pre-loading and refresh of thousands of certificates and
large certificate revocation lists. Privacy comes only from changing identifiers, and the authority
stays on the trust path. Zhong et al.~\cite{zhong2019aggregation} compress many signatures into one
with certificateless aggregate signatures, which cuts verification and communication, but the trace
authority remains a single point of trust and no ledger shares revocation state between RSUs.

\section{ECC and Certificateless Authentication}

Liu et al.~\cite{liu2025ecc} propose cooperative ECC-based mutual authentication with pseudonyms,
batch verification and group keys for low RSU workload. The trusted authority issues and holds every
long-term key, so it can link any pseudonym to a real identity. Qi et al.~\cite{qi2022certificateless}
combine pseudonyms with certificateless signatures to avoid certificate management and key escrow.
Certificateless cryptography removes the certificate, not the authority: a key generation centre
still issues partial keys, and per-message verification remains.

\section{Hardware and Game-Theoretic Approaches}

Li et al.~\cite{li2024puf} derive authentication secrets from physical unclonable functions (PUFs),
which resist key extraction but require PUF hardware in every OBU and cannot be retrofitted to a mixed
fleet. Gupta et al.~\cite{gupta2022game} combine a game-theoretic trust model with a blockchain to
decide authentication, at the cost of decision and consensus complexity.

\section{Blockchain-Based Authentication}

Blockchains give RSUs shared, tamper-evident state. BAAIoV~\cite{alsaeed2025baaiov} uses a blockchain
for both authentication and authorization in the Internet of Vehicles, and BPA~\cite{li2024bpa} builds
privacy-preserving authentication on it. B-TSCA~\cite{wang2021btsca} uses blockchain-assisted
trustworthiness computation so that re-authentication at the next RSU is cheaper than the first
authentication. Xie et al.~\cite{xie2023handover} add PUFs for lightweight V2I handover and V2V
broadcast authentication. In all of these the ledger supplies shared state, but the credential is
still shown to the verifier, so privacy depends on pseudonym management and nothing prevents linking
two sessions of the same vehicle.

\section{Zero-Knowledge Authentication}

Rasheed et al.~\cite{rasheed2020adaptive} use group-based zero-knowledge proofs with configurable
privacy and revocation. They rely on interactive proof rounds and a trusted group manager, which costs
round trips at every handover. ZAMA~\cite{xi2022zama} combines a ZKP, ECC and a Fujisaki--Okamoto
commitment for anonymous mutual authentication with traceability, at a higher cost per
authentication and with a central authority. Varma and Kumar~\cite{varma2023zkp} focus on lightweight
ZKP authentication during handover. Without decentralised state, authorization and revocation
decisions cannot follow a vehicle between RSUs and cannot be audited independently.

\section{Blockchain with Zero-Knowledge Proofs}

The schemes closest to this work combine both. Jiang and Guo~\cite{jiang2024trug} use master--slave
chains with a PBFT variant and a ZKP. Wang et al.~\cite{wang2025privacy} use identity-based polynomial
commitments, a ZKP and a Merkle Patricia trie on a blockchain for anonymous authentication and data
retrieval; a key generation centre remains in the setup, and service-level policy, the revocation
lifecycle and RSU handover are not modelled. Li et al.~\cite{li2026emergency} authenticate emergency
vehicles with one-time hash commitments, category Merkle trees, a ZKP and on-chain state, but cover a
small static fleet under one authorization decision.

\section{Misbehaviour Detection and Location Privacy}

Plausibility checks and learned detectors are evaluated on VeReMi~\cite{heijden2018veremi} and
VeReMi Extension~\cite{kamel2020veremiext}. For privacy against a global eavesdropper, pseudonym
changes alone are weak because a tracker links the old and the new pseudonym by position and
speed~\cite{gruteser2005tracking}; mix zones~\cite{beresford2003mix,freudiger2007mix}, where vehicles
change pseudonyms together and stay silent briefly, break that link.

\section{Comparison and Research Gap}

\begin{table}[htbp]\centering\small
\caption{Families of VANET authentication schemes (adapted from the Phase~I review)}\label{tab:families}
\resizebox{\textwidth}{!}{%
\begin{tabular}{@{}lllll@{}}
\toprule
Family & Identity privacy & Shared state & Accountability & Main limitation \\
\midrule
PKI / pseudonym~\cite{raya2007securing} & medium--high & low & high & certificates and revocation lists \\
ECC / certificateless~\cite{liu2025ecc,qi2022certificateless} & medium & low & high & dependence on the TA \\
Aggregate~\cite{zhong2019aggregation} & high & low & medium--high & trace authority \\
PUF~\cite{li2024puf} & high & low & medium & special hardware \\
Blockchain~\cite{alsaeed2025baaiov,li2024bpa,wang2021btsca,xie2023handover} & medium & high & high & consensus latency \\
ZKP~\cite{varma2023zkp,rasheed2020adaptive,xi2022zama} & very high & low & medium--high & proof cost, no shared state \\
Blockchain + ZKP~\cite{jiang2024trug,wang2025privacy,li2026emergency} & very high & high & high & integration, scalability \\
\textbf{This work} & very high & high & high & measured in Chapter~\ref{ch:results} \\
\bottomrule
\end{tabular}}
\end{table}

Table~\ref{tab:families} summarises the families. The direct blockchain-plus-ZKP literature shows
that the combination is feasible. What it does not provide is a complete lifecycle (handover,
renewal, misbehaviour reporting, revocation that reaches vehicles, threshold identity reveal,
per-service policy) evaluated on a realistic network with measured costs. This project therefore
does not claim the combination as new. Its contribution is a concrete integration that keeps the
ledger off the safety-message path, and its measurement against certificate-based and centralised
baselines on ns-3 with real traffic.
